\clearpage
\fancyhead[L]{\small HUMBOLDT. CONSPECTUS LONGITUDINUM ETC.}
\hypertarget{native-514-lower}{}\label{native-514-lower}
\begin{center}\small Göttingische gelehrte Anzeigen. Stück 193. S. 1928. 3. Dezember 1808.\\[3pt]\rule{.72\linewidth}{.3pt}\end{center}

\emph{Conspectus longitudinum et latitudinum geographicarum, per decursum annorum} 1799 \emph{ad} 1804 \emph{in plaga aequinoctiali ab} \textsc{Alex. de Humboldt} \emph{astronomice observatarum, calculo subjecit} \textsc{Jabbo Oltmans}. 16 Seiten im Folioformat. Bei \textsc{F. Schöll} und \textsc{J. G. Cotta} in Paris und Tübingen.

Den wertvollen Ergebnissen, die die Geografie der amerikanischen Tropenländer der Reise von Herrn \textsc{von Humboldt} verdankt, wird bekanntlich ein eigener Teil seines großen Werks gewidmet. Darin dürfen wir alle seine zahlreichen astronomischen Beobachtungen in voller Länge erwarten, dazu eine sehr sorgfältige Auswertung dieser Beobachtungen durch den bereits angesehenen Herrn \textsc{Jabbo Oltmans}. Da der Druck dieses Werks jedoch beträchtliche Zeit beanspruchen wird, hat man inzwischen diese Übersicht der wichtigsten Ortsbestimmungen zusammengestellt. Sie enthält nur die Endergebnisse. Fast dreihundert Bestimmungen werden hier mitgeteilt, die größtenteils auf den eigenen Beobachtungen von Herrn \textsc{von Humboldt} beruhen. Man kann nicht umhin, die Tatkraft und Vielseitigkeit dieses Mannes zu bewundern. Er verfolgte so viele verschiedenartige Ziele zugleich und wusste doch jedes so zu erreichen, als hätte er ihm seine Zeit und Kräfte ungeteilt widmen können.
